\section{Measurements}

During the final part of the course, electrical measurements were performed on the process control module (PCM) structures of the BiCMOS5 wafer: Van der Pauw (VdP) squares, electrical line width measurement (ELM) resistors, and NMOS/PMOS transistors in all four $V_\mathrm{T}$-adjust quadrants.

The measurements were taken on a Cascade 31 automatic probe station and a Keysight B1500A parameter analyzer controlled by IC-CAP, using four-point probing throughout so that contact and needle resistances drop out.

Each wafer quadrant has a different boron $V_\mathrm{T}$-adjust dose: $0$, $3\times10^{11}$, $6\times10^{11}$ and $9\times10^{11}\,\mathrm{cm}^{-2}$. Junction depths $x_\mathrm{j}$ for the N-Well, SN and SP implants were taken from the Sentaurus simulations in Chapter~3, and the in-line sheet resistances from Chapter~4, so the measurements here can be directly compared with the earlier process and simulation results.

\subsection{Van der Pauw Sheet Resistance and Doping Concentration}

The VdP structures were measured in the top-left quadrant ($V_\mathrm{T}$-adjust dose $0$) using a symmetric four-point configuration. The analyzer swept the force voltage between $-0.2$\,V and $+0.2$\,V at grounded substrate, and the sheet resistance $R_\mathrm{s}$ was obtained from the least-squares slope of $\Delta V$ versus $I_\mathrm{force}$ followed by the Van der Pauw correction factor
\[
  R_\mathrm{s} = \frac{\pi}{\ln 2} R_{AB,CD} \approx 4.532\,R_{AB,CD}.
\]
Within this $\pm 0.2$\,V window the pn-junctions to the substrate remain strongly reverse-biased, so the conduction is confined to the implanted layer and the VdP assumptions hold.

The measured resistances and fitted sheet resistances are summarized in Table~\ref{tab:ch7-vdp-rsheet}. The last column shows the in-line monitor values from Chapter~4 for reference.

\begin{table}[htbp]
  \centering
  \caption{Van der Pauw sheet resistances of the implanted layers and metal, compared to the in-line monitor values from Chapter~4.}
  \label{tab:ch7-vdp-rsheet}
  \begin{tabular}{lrrrr}
    \toprule
    Layer & $R_{AB,CD}$ ($\Omega$) & $R_\mathrm{s}$ ($\Omega/\square$) & Fit $R^2$ & Ch.~4 $R_\mathrm{s}$ ($\Omega/\square$) \\
    \midrule
    N-Well      & 384.02  & 1740.54 & 0.9991 & 1488.8 \\
    SN          &  13.28  &   60.19 & 0.9991 &   57.4 \\
    SP (in N-Well) & 107.75 & 488.37 & 0.9990 & 510.6 \\
    IC metal    &   0.0352 & 0.1597 & 0.9939 & -- \\
    \bottomrule
  \end{tabular}
\end{table}

The SN and SP values agree with the Chapter~4 monitor resistances to within about $5\%$, while the N-Well shows a larger deviation of roughly $17\%$. This is consistent with the N-Well being the most lightly doped layer, and therefore most sensitive to wafer position, surface depletion, and differences between blanket monitor wafers and the drawn VdP cloverleaf on the PCM.

The IC metal has a sheet resistance of only $0.1597\,\Omega/\square$, which is three to four orders of magnitude lower than the implanted silicon layers, whose values range from about $60\,\Omega/\square$ for SN to $1740\,\Omega/\square$ for the N-Well. This large difference is expected because the IC layer is a $0.2\,\mu\mathrm{m}$ thick Al--1\%Si film, whereas the NW, SN and SP layers conduct through doped silicon. Metals contain a much higher free-carrier density than doped semiconductor layers, so their resistivity is much lower even when the carrier mobility in lightly doped silicon is relatively high. For this reason, the metal layer is used for low-resistance interconnects, while the implanted layers are suitable for forming resistive structures on the chip.

Using the simulated junction depths $x_\mathrm{j}$ from Chapter~3, the depth-averaged resistivity of each implanted layer follows
\[
  \rho = R_\mathrm{s}\,x_\mathrm{j}.
\]
The average active dopant concentration $\bar N$ is then obtained from the Irvin curves for the corresponding carrier type. The resulting resistivities and concentrations are given in Table~\ref{tab:ch7-vdp-conc}.

\begin{table}[htbp]
  \centering
  \caption{Depth-averaged resistivity and active dopant concentration from the VdP measurements and Irvin curves.}
  \label{tab:ch7-vdp-conc}
  \begin{tabular}{lcccc}
    \toprule
    Layer & $x_\mathrm{j}$ ($\mu$m) & $\rho = R_\mathrm{s}x_\mathrm{j}$ ($\Omega\cdot\mathrm{cm}$) & Type & $\bar N$ ($\mathrm{cm}^{-3}$) \\
    \midrule
    N-Well & 2.085 & 0.363            & n & $1.5\times10^{16}$ \\
    SN     & 0.537 & $3.2\times10^{-3}$ & n & $2.3\times10^{19}$ \\
    SP     & 0.594 & $2.9\times10^{-2}$ & p & $1.6\times10^{18}$ \\
    \bottomrule
  \end{tabular}
\end{table}

These concentrations match the Chapter~4 Irvin-curve values within about $10\%$ for all three implants, and they sit below the simulated peak concentrations from Chapter~3 because they represent averages over the diffused junction profiles rather than the peak implant doses.
\subsection{Electrical Line Width Measurements and Lateral Out-Diffusion}

To probe lateral out-diffusion of dopants, ELM structures with widths $W = 2$, $5$ and $10\,\mu\mathrm{m}$ and length $L = 206\,\mu\mathrm{m}$ were measured for the N-Well and SN implants. For each resistor, a four-point measurement was taken and the apparent sheet resistance was computed as
\[
  R_\mathrm{s,app} = R \frac{W}{L},
\]
using the drawn width $W$. The results are listed in Table~\ref{tab:ch7-elm-raw}.

\begin{table}[htbp]
  \centering
  \caption{Raw ELM resistances and apparent sheet resistances for N-Well and SN resistors ($L = 206\,\mu\mathrm{m}$).}
  \label{tab:ch7-elm-raw}
  \begin{tabular}{lrrr}
    \toprule
    Layer & $W$ ($\mu$m) & $R$ ($\Omega$) & $R_\mathrm{s,app}$ ($\Omega/\square$) \\
    \midrule
    N-Well &  2 & 58832.7 & 571.2 \\
    N-Well &  5 & 38478.1 & 933.9 \\
    N-Well & 10 & 19689.2 & 955.8 \\
    SN     &  2 &  5244.8 &  50.9 \\
    SN     &  5 &  2273.5 &  55.2 \\
    SN     & 10 &  1171.1 &  56.8 \\
    \bottomrule
  \end{tabular}
\end{table}

For both implants the apparent sheet resistance is systematically lower than the VdP value when the smallest width is used, and the discrepancy increases as the drawn line becomes narrower. This indicates that the conducting strip is wider than the nominal geometry due to lateral dopant diffusion, so the effective width is
\[
  W_\mathrm{eff} = W + 2\Delta W.
\]

Assuming a constant out-diffusion $\Delta W$ per side, straight-line fits of $R$ versus $1/(W + 2\Delta W)$ yield both the true sheet resistance and the lateral out-diffusion. The fitted values are summarized in Table~\ref{tab:ch7-elm-fit}.

\begin{table}[htbp]
  \centering
  \caption{Sheet resistance and lateral out-diffusion extracted from the ELM structures.}
  \label{tab:ch7-elm-fit}
  \begin{tabular}{lrrr}
    \toprule
    Layer & $R_\mathrm{s,ELM}$ ($\Omega/\square$) & $\Delta W$ ($\mu$m) & Fit $R^2$ \\
    \midrule
    N-Well & 1129.2 & 0.80 & 0.986 \\
    SN     &   58.6 & 0.15 & 0.999999 \\
    \bottomrule
  \end{tabular}
\end{table}

For the SN implant the ELM sheet resistance agrees with the VdP value ($60.2\,\Omega/\square$) to within about $3\%$, and the extracted lateral out-diffusion is only $\Delta W \approx 0.15\,\mu\mathrm{m}$ per side. The N-Well lines show a much larger effective widening: fitting the ELM data gives $\Delta W \approx 0.8\,\mu\mathrm{m}$, while anchoring the sheet resistance to the VdP value suggests an out-diffusion closer to $2.1\,\mu\mathrm{m}$ per side, comparable to the vertical junction depth.

This difference is explained by both the longer thermal budget and the faster diffusivity of phosphorus in the N-Well compared to arsenic in the SN: the N-Well receives a $4$\,h, $1150^\circ$C drive-in, so the dopants spread several micrometres sideways as well as vertically, whereas the SN sees only the shorter gate oxidation anneal.

\subsection{Transistor Characterisation: Threshold Voltage, Mobility and Short-Channel Effects}

NMOS and PMOS transistors with gate geometries $W:L = 20\!:\!1$, $20\!:\!2$, $20\!:\!5$, $20\!:\!10$ and $60\!:\!2$ were measured in the bottom-right quadrant ($V_\mathrm{T}$-adjust dose $9\times10^{11}\,\mathrm{cm}^{-2}$). Transfer characteristics $I_\mathrm{D}(V_\mathrm{G})$ were recorded at $|V_\mathrm{D}| = 0.1$\,V, and threshold voltages and mobilities were extracted using the long-channel unified model.

In the linear region ($|V_\mathrm{D}|$ small), the drain current follows
\[
  I_\mathrm{D} = \mu \frac{W}{L} C_\mathrm{ox} \left( V_\mathrm{GS} - V_\mathrm{T} - \frac{1}{2}V_\mathrm{DS} \right) V_\mathrm{DS},
\]
so $I_\mathrm{D}$ is linear in $V_\mathrm{D}$ and proportional to $W/L$. Threshold voltage $V_\mathrm{T}$ was extracted by linear extrapolation of $I_\mathrm{D}(V_\mathrm{G})$ at the point of maximum transconductance $g_\mathrm{m,max}$, and the mobility $\mu$ was obtained from
\[
  \mu = \frac{g_\mathrm{m,max}L}{C_\mathrm{ox}WV_\mathrm{DS}}.
\]

The extracted values for the $9\times10^{11}\,\mathrm{cm}^{-2}$ quadrant are listed in Table~\ref{tab:ch7-transfers}.

\begin{table}[htbp]
  \centering
  \caption{Threshold voltages and mobilities extracted from the transfer characteristics in the $9\times10^{11}\,\mathrm{cm}^{-2}$ quadrant ($V_\mathrm{D} = \pm 0.1$\,V).}
  \label{tab:ch7-transfers}
  \begin{tabular}{lrrrr}
    \toprule
    Type & $W:L$ & $V_\mathrm{T}$ (V) & $g_\mathrm{m,max}$ ($\mu$S) & $\mu$ ($\mathrm{cm}^2/\mathrm{V\,s}$) \\
    \midrule
    NMOS & 20:1  &  0.572 & 68.0  & 985 \\
    NMOS & 20:2  &  0.740 & 27.5  & 796 \\
    NMOS & 20:5  &  0.819 & 10.1  & 734 \\
    NMOS & 20:10 &  0.811 & 5.18  & 750 \\
    NMOS & 60:2  &  0.779 & 80.6  & 778 \\
    \midrule
    PMOS & 20:1  & -1.733 & 39.0  & 565 \\
    PMOS & 20:2  & -3.262 & 10.4  & 301 \\
    PMOS & 20:5  & -3.300 & 3.25  & 236 \\
    PMOS & 20:10 & -3.305 & 1.59  & 230 \\
    PMOS & 60:2  & -3.213 & 31.5  & 304 \\
    \bottomrule
  \end{tabular}
\end{table}

For gate lengths $L \ge 2\,\mu\mathrm{m}$, both the threshold voltage and the mobility remain nearly constant as geometry changes, confirming that these are process parameters rather than device parameters. The large-area devices show electron mobilities around $750$--$800\,\mathrm{cm}^2/\mathrm{V\,s}$ and hole mobilities around $230$--$300\,\mathrm{cm}^2/\mathrm{V\,s}$, roughly a factor of three difference as expected from bulk silicon.

At the shortest gate length $L = 1\,\mu\mathrm{m}$ (devices $20:1$), clear short-channel effects appear. The NMOS threshold voltage rolls off from about $0.82$\,V to $0.57$\,V, and the extracted mobility is artificially inflated, reflecting effective channel shortening and drain-induced barrier lowering rather than a true change in carrier transport. The PMOS devices show an even stronger threshold reduction, from about $-3.30$\,V to $-1.73$\,V, consistent with deeper junctions and the lightly doped N-Well underneath.

\subsubsection{Threshold Voltage versus $V_\mathrm{T}$-Adjust Dose}

To quantify the effect of the boron $V_\mathrm{T}$-adjust implant, the $W:L = 20:5$ transistors were measured in all four quadrants with doses $0$, $3\times10^{11}$, $6\times10^{11}$ and $9\times10^{11}\,\mathrm{cm}^{-2}$. The extracted threshold voltages are listed in Table~\ref{tab:ch7-vt-dose}.

\begin{table}[htbp]
  \centering
  \caption{Threshold voltage versus $V_\mathrm{T}$-adjust dose for NMOS and PMOS transistors ($W:L = 20:5$).}
  \label{tab:ch7-vt-dose}
  \begin{tabular}{lrr}
    \toprule
    Dose ($\mathrm{cm}^{-2}$) & $V_\mathrm{T,\ NMOS}$ (V) & $V_\mathrm{T,\ PMOS}$ (V) \\
    \midrule
    $0\times10^{11}$ & -0.034 & -3.702 \\
    $3\times10^{11}$ &  0.388 & -3.511 \\
    $6\times10^{11}$ &  0.432 & -3.413 \\
    $9\times10^{11}$ &  0.819 & -3.300 \\
    \bottomrule
  \end{tabular}
\end{table}

The NMOS threshold voltage increases monotonically from near zero to about $0.82$\,V as the boron dose increases, because the additional acceptor charge in the channel depletion region must be balanced by a larger gate voltage. In contrast, the same boron implant partially compensates the N-Well underneath the PMOS, making its threshold voltage less negative by roughly $0.4$\,V over the same dose range.

This behaviour matches the qualitative theory from Chapter~2 and the Sentaurus device simulations from Chapter~3: boron $V_\mathrm{T}$-adjust raises NMOS $V_\mathrm{T}$ while making PMOS easier to turn on.

\subsubsection{Output Characteristics and Channel-Length Modulation}

Output characteristics $I_\mathrm{D}(V_\mathrm{D})$ were measured for the $W:L = 20:5$ and $20:1$ devices at the highest $V_\mathrm{T}$-adjust dose, with gate voltage steps spanning the linear and saturation regions. In long channels, the saturation current ideally depends only on the gate overdrive and is flat versus $V_\mathrm{D}$, but real devices show finite slopes due to channel-length modulation and velocity saturation.

Channel-length modulation parameters $\lambda$ were extracted from linear fits of $I_\mathrm{D}$ versus $V_\mathrm{D}$ in the saturation region, where $|V_\mathrm{D}|$ exceeds $|V_\mathrm{G} - V_\mathrm{T}| + 0.3$\,V. Representative values are shown in Table~\ref{tab:ch7-lambda}.

\begin{table}[htbp]
  \centering
  \caption{Representative channel-length modulation parameters for NMOS and PMOS transistors.}
  \label{tab:ch7-lambda}
  \begin{tabular}{lrrr}
    \toprule
    Type & $W:L$ & $V_\mathrm{G}$ (V) & $\lambda$ ($\mathrm{V}^{-1}$) \\
    \midrule
    NMOS & 20:5 & 1--5   & $\approx 0.022$ \\
    NMOS & 20:1 & 2--4   & $\approx 0.16$ \\
    PMOS & 20:5 & -4 to -7 & $\approx 0.016$ \\
    PMOS & 20:1 & various  & $\approx 0.1$--1 (no true saturation) \\
    \bottomrule
  \end{tabular}
\end{table}

The long-channel devices ($W:L = 20:5$) show small $\lambda$ values around $0.02\,\mathrm{V}^{-1}$, indicating weak channel-length modulation and clean saturation. For the short-channel $20:1$ devices, $\lambda$ increases by almost an order of magnitude, and the PMOS $20:1$ transistor never fully saturates within $|V_\mathrm{D}| \le 5$\,V, consistent with drain-induced barrier lowering and punch-through in the deep, lightly doped N-Well.

These measurements reproduce the qualitative trends seen in the Sentaurus device simulations in Chapter~3: long-channel transistors follow the square-law model with modest saturation slopes, while short-channel devices exhibit threshold roll-off, increased apparent mobility, larger $\lambda$ and signatures of velocity saturation.